\section{Introduzione}
\label{sec:introduzione}

\subsection{Contesto: Semantic Web e Ontologie}

Il Semantic Web, introdotto da Tim Berners-Lee nel 2001, rappresenta un'estensione
del Web tradizionale in cui l'informazione \`e strutturata e dotata di semantica
formale, consentendo alle macchine di interpretare e ragionare sui dati in modo
automatico. Il fondamento teorico di questa visione risiede nel concetto di
\textbf{ontologia}, definita da \citet{gruber1993} come \textless{}una specifica
esplicita e formale di una concettualizzazione condivisa\textgreater{}.

Le ontologie, espresse in linguaggi formali come OWL 2 \citep{w3c-owl2-primer},
permettono di modellare un dominio di conoscenza attraverso:
\begin{itemize}[nosep]
  \item \textbf{Classi} --- categorie di entit\`a (es. \texttt{VideoGame}, \texttt{Developer});
  \item \textbf{Propriet\`a} --- relazioni tra entit\`a (es. \texttt{developedBy}) e attributi (es. \texttt{metacriticScore});
  \item \textbf{Assiomi} --- vincoli logici che arricchiscono la semantica (es. classi definite, disjointness, catene di propriet\`a).
\end{itemize}

Un knowledge graph \`e la concretizzazione pratica di un'ontologia popolata con
dati reali: nodi (istanze) e archi (asserzioni RDF) formano un grafo interrogabile
tramite query SPARQL.

\subsection{Obiettivo del Progetto}

Il progetto \textbf{Videogame Semantic Search} si propone di:
\begin{enumerate}[nosep]
  \item \textbf{Costruire un'ontologia OWL 2 di alta qualit\`a} sul dominio dei videogiochi, coprendo il periodo 2010--2026, con uno schema ricco di assiomi (classi primitive e definite, disjointness, property chain, allineamenti esterni);
  \item \textbf{Popolare l'ontologia da Wikidata} tramite query SPARQL all'endpoint pubblico di Wikidata, producendo un insieme di istanze di oltre 1.1 milioni di triple;
  \item \textbf{Abilitare l'interrogazione in linguaggio naturale} attraverso un agente SPARQL basato su Large Language Model (GPT-4.1-mini) che traduce domande in italiano (o inglese) in query SPARQL eseguibili;
  \item \textbf{Visualizzare i risultati come grafo di conoscenza interattivo}, navigabile con zoom, pan, filtri e dettaglio nodi.
\end{enumerate}

\subsection{Casi d'Uso}

I principali casi d'uso del sistema sono:
\begin{itemize}[nosep]
  \item \textbf{Ricerca semantica} --- interrogare il knowledge graph con domande come \textless{}Quali giochi ha sviluppato FromSoftware?\textgreater{} o \textless{}Top 10 giochi RPG con Metacritic pi\`u alto\textgreater{};
  \item \textbf{Grafo di conoscenza interattivo} --- esplorare visualmente le relazioni tra giochi, sviluppatori, generi, piattaforme e franchise;
  \item \textbf{Reasoning automatico} --- classificare automaticamente i giochi in categorie derivate (MultiplayerGame, HighlyRatedGame, SequelGame, ecc.) sfruttando gli assiomi OWL 2.
\end{itemize}

\subsection{Panoramica della Relazione}

Il presente documento \`e organizzato come segue.
Le Sezioni~\ref{sec:ontologia_overview}--\ref{sec:profilo_dl} costituiscono il cuore
della relazione e analizzano in dettaglio l'ontologia: panoramica, classi, propriet\`a
e chiavi, classi definite, disjointness, allineamento esterno, e infine espressivit\`a
DL e profilo OWL 2.
La Sezione~\ref{sec:reasoning} approfondisce il sistema di reasoning a due livelli con
un esempio completo di inferenza, incluso il problema degli \texttt{owl:sameAs} falsi e
la sua risoluzione; la Sezione~\ref{sec:interrogazione} mostra come interrogare con
SPARQL la conoscenza inferita.
Le Sezioni~\ref{sec:popolamento}, \ref{sec:architettura} e \ref{sec:frontend_backend}
descrivono il contesto: la pipeline di popolamento da Wikidata, l'architettura software
e l'applicazione web.
La Sezione~\ref{sec:conclusione} riassume i risultati e delinea possibili sviluppi futuri.
Le appendici riportano le istanze di test, la sintassi DL delle classi definite e gli
assiomi principali dello schema in Turtle.
